
%%%%%%%%%%%%%%%%%%%%%%%%%%%%%%%%%%%%%%%%
%           main body
%%%%%%%%%%%%%%%%%%%%%%%%%%%%%%%%%%%%%%%%

%-----------------------------------
%	TITLE PAGE
%-----------------------------------

\title[数学分析]{\texorpdfstring{\LARGE \bfseries 第九章\quad 数项级数}{第九章 数项级数}} % The short title appears at the bottom of every slide, the full title is only on the title page
\author[正项级数] % Your institution as it will appear on the bottom of every slide, maybe shorthand to save space
{\Large \bfseries 正项级数}
% \author{Singular Value Decomposition} % Your name
% \date{\today} % Date, can be changed to a custom date
\date{}

%------------------------------------------------

\begin{document}

%------------------------------------------------

\begin{frame}

\titlepage % Print the title page as the first slide

\end{frame}

%------------------------------------------------

% \begin{frame}
% \frametitle{内容提要} % Table of contents slide, comment this block out to remove it
% \tableofcontents % Throughout your presentation, if you choose to use \section{} and \subsection{} commands, these will automatically be printed on this slide as an overview of your presentation
% \end{frame}

%-----------------------------------
%	PRESENTATION SLIDES
%-----------------------------------

%------------------------------------------------

\begin{frame}{引言}

\begin{columns}
\begin{column}{0.51\textwidth}
\centering{
\includegraphics[width=0.75\textwidth]{images/zeno-paradox.jpg}
}
\end{column}
\begin{column}{0.45\textwidth}
$S_1(1 + q + q^2 + \cdots + q^n + \cdots)$
\end{column}
\end{columns}

\vspace*{0.3em}

\begin{columns}
\begin{column}{0.51\textwidth}
\centering{

\includegraphics[width=0.75\textwidth]{images/一尺之捶-日取其半-万世不竭.jpg}
}
\end{column}
\begin{column}{0.45\textwidth}
$\dfrac{1}{2} + \dfrac{1}{4} + \cdots + \dfrac{1}{2^n} + \cdots$
\end{column}
\end{columns}

\pause
\vspace*{0.3em}

\begin{center}
这些级数有共同的特点: 通项都是正数!
\end{center}

\end{frame}

%------------------------------------------------

\section[收敛原理]{正项级数的收敛原理}

%------------------------------------------------

\begin{frame}{正项级数的定义}

\begin{Def}
如果级数 $\sum\limits_{n=1}^\infty x_n$ 的通项都是非负实数, 即
\[x_n \geqslant 0, \quad n = 1, 2, \cdots,\]
则称此级数为\textbf{正项级数}.
\end{Def}

\pause

显然, 正项级数 $\sum\limits_{n=1}^\infty x_n$ 的部分和数列 $\{S_n\} $是单调增加的, 即
\[S_n = \sum\limits_{k=1}^n x_k \leqslant \sum\limits_{k=1}^{n+1} x_k = S_{n+1}, \quad n = 1, 2, \cdots.\]
由数列的单调有界收敛定理, 立刻可以得到

\end{frame}

%------------------------------------------------

\begin{frame}{正项级数的收敛原理}

\begin{thm}[正项级数的收敛原理]
\begin{center}
正项级数收敛的充分必要条件是它的部分和数列有上界.
\end{center}
\end{thm}

\pause
\vspace*{0.4em}

\textbf{例:} 考虑级数 $\sum\limits_{n=2}^\infty \left[ \dfrac{1}{\nroot[n]{n}} \ln{\dfrac{n^2}{(n-1)(n+1)}} \right].$ 容易看出它是一个正项级数.
\pause
它的部分和数列的通项为
\begin{align*}
S_n & = \sum\limits_{k=2}^{n+1} \left[ \dfrac{1}{\nroot[k]{k}} \ln{\dfrac{k^2}{(k-1)(k+1)}} \right] ~ {\color{red}<} \sum\limits_{k=2}^{n+1} ~ {\color{red}\cancel{\dfrac{1}{\nroot[k]{k}}}} \ln{\left( \left.\dfrac{k}{k-1} \right/ \dfrac{k+1}{k} \right)} \\
& = \sum\limits_{k=2}^{n+1} \left( \ln{\dfrac{k}{k-1}} - \ln{\dfrac{k+1}{k}} \right) = \ln{2} - \ln{\dfrac{n+2}{n+1}} < \ln{2},
\end{align*}
所以正项级数 $\sum\limits_{n=2}^\infty \left[ \dfrac{1}{\nroot[n]{n}} \ln{\dfrac{n^2}{(n-1)(n+1)}} \right]$ 收敛.

\end{frame}

%------------------------------------------------

\begin{frame}{正项级数的收敛原理}

在这个例子中, 我们实际上使用了马上要介绍的“比较判别法”, 即使用了不等式
\[\dfrac{1}{\nroot[n]{n}} \ln{\dfrac{n^2}{(n-1)(n+1)}} := x_n \leqslant y_n := \ln{\dfrac{n^2}{(n-1)(n+1)}},\]
$\sum\limits_{n=2}^\infty y_n$ 相对比较好求, 敛散性相对比较好判断. 可以通过判断它的收敛性, 进而得出 $\sum\limits_{n=2}^\infty x_n$ 的收敛性.

\pause
\vspace*{1em}

本节剩下的内容, 就是介绍一系列包括“比较判别法”在内的简便实用的判别正项级数敛散性的方法.

\end{frame}

%------------------------------------------------

\section[比较判别法]{正项级数敛散性的比较判别法}

%------------------------------------------------

\begin{frame}{比较判别法}

判断一个正项级数是否收敛, 最常用的方法是用一个已知敛散性的级数与它进行比较.

\pause

\begin{thm}[比较判别法]
设 $\sum\limits_{n=1}^\infty x_n$ 与 $\sum\limits_{n=1}^\infty y_n$ 是两个正项级数, 若存在常数 $A > 0,$ 使得
\[x_n \leqslant Ay_n, \quad n = 1, 2 ,\cdots,\]
\begin{enumerate}
\item[(1)] 当 $\sum\limits_{n=1}^\infty y_n$ 收敛时, $\sum\limits_{n=1}^\infty x_n$ 也收敛;
\item[(2)] 当 $\sum\limits_{n=1}^\infty x_n$ 发散时, $\sum\limits_{n=1}^\infty y_n$ 也发散.
\end{enumerate}
\end{thm}

\pause

\begin{empheq}[box={\fancymath[colback=cyan!30, drop lifted shadow, sharp corners]}]{equation*}
\text{\bfseries\color{red!80} 大的收敛, 则小的也收敛; \quad 小的发散, 则大的也发散.}
\end{empheq}

\end{frame}

%------------------------------------------------

\begin{frame}{比较判别法的证明}

设级数 $\sum\limits_{n=1}^\infty x_n$ 与 $\sum\limits_{n=1}^\infty y_n$ 的部分和序列分别为 $\{S_n\}$ 与 $\{T_n\},$ 那么显然有
\[S_n \leqslant AT_n, \quad n = 1, 2 ,\cdots,\]
\pause
于是当数列 $\{T_n\}$ 有上界时, $\{S_n\}$ 也有上界; 当 $\{S_n\}$ 无上界时, $\{T_n\}$ 必定也无上界. 那么由~{\color{red}正项级数的收敛原理}\tikzmark{mark1}

\pause

\begin{tikzpicture}[remember picture, overlay]
\node[rectangle, draw=black, fill=cyan!30, text width=0.35\textwidth] at ([xshift=230, yshift=-25]mark1) (mark1_to) {{\bfseries\smaller[2]正项级数收敛的~{\color{red}充分必要条件}~\color{themecolor}是它的部分和数列有上界}};
\draw[->, color=black] (mark1_to.north west) -- (mark1);
\end{tikzpicture}

\pause

\begin{enumerate}
\item[(1)] 当 $\sum\limits_{n=1}^\infty y_n$ 收敛时, $\sum\limits_{n=1}^\infty x_n$ 也收敛;
\item[(2)] 当 $\sum\limits_{n=1}^\infty x_n$ 发散时, $\sum\limits_{n=1}^\infty y_n$ 也发散.
\end{enumerate}

\pause

\begin{attnblock}
由于改变级数有限个项的数值, 并不会改变它的敛散性, 所以比较判别法的条件可以放宽为: “存在正整数 $N$ 与常数 $A>0,$ 使得 $x_n \leqslant Ay_n$ 对一切 $n > N$ 成立”.
\end{attnblock}

\end{frame}

%------------------------------------------------

\begin{frame}{比较判别法的应用示例}

\begin{eg}
\begin{center}
判断正项级数 $\sum\limits_{n=1}^\infty \dfrac{n+3}{2n^3-n}$ 的敛散性.
\end{center}
\end{eg}

\pause

\startproof[bg=yellow, fg=red](解)
容易看出, 当 $n > 3$ 时成立
\[\dfrac{n + ~{\color{red}3}}{2n^3 - ~ {\color{red}n}} < \dfrac{n+{\color{red}n}}{2n^3 - ~{\color{red}n^3}} = \dfrac{2}{n^2},\]
由 $\sum\limits_{n=1}^\infty \dfrac{1}{n^2}$ 的收敛性 (值等于 $\dfrac{\pi^2}{6}$), 知 $\sum\limits_{n=1}^\infty \dfrac{n+3}{2n^3-n}$ 收敛.
\qed

\end{frame}

%------------------------------------------------

\begin{frame}{比较判别法的应用示例}

\begin{eg}
\begin{center}
判断正项级数 $\sum\limits_{n=1}^\infty \sin\dfrac{\pi}{n}$ 的敛散性.
\end{center}
\end{eg}

\pause

\begin{columns}
\begin{column}{0.4\textwidth}
\centering{
\scalebox{1.1}{
\begin{tikzpicture}[scale=1.5, domain=0:3.1416/2]
\draw[->] (-0.3,0) -- (2.2,0) node[right] {$x$};
\draw[->] (0,-0.3) -- (0,2.2) node[above] {$f(x)$};

\draw[color=orange,dashed,thick,domain=0:pi/2,samples=100]
plot (\x,\x) node[above right] {$f(x)=x$};

\draw[color=blue,thick,domain=0:pi/2,samples=100]
plot (\x,{sin(\x r)}) node[above right] {$f(x)=\sin x$};

\draw[color=red,dashed,thick,domain=0:pi/2,samples=100]
plot (\x,{2*\x/pi}) node[below right] {$f(x)=\dfrac{2}{\pi}x$};

\draw[color=themecolor,dashed]
(pi/2,0) node[below] {\smaller[1]$\dfrac{\pi}{2}$} -- (pi/2,pi/2);

\draw[color=themecolor,dashed]
(0,pi/2) node[left] {\smaller[1]$\dfrac{\pi}{2}$} -- (pi/2,pi/2);

\draw[color=themecolor,dashed]
(0,1) node[left] {\smaller[1]$1$} -- (pi/2,1);
\end{tikzpicture}
}
}
\end{column}
\begin{column}{0.58\textwidth}
从右图可看出, 当 $x \in [0, \dfrac{\pi}{2}]$ 时, 成立不等式 $\sin x \geqslant \dfrac{2}{\pi} x,$
\pause
所以当 $n\geqslant 2$ 时有
\[\sin\dfrac{\pi}{n} \geqslant \dfrac{2}{\pi} \cdot \dfrac{\pi}{n} = \dfrac{2}{n}.\]
由于 $\sum\limits_{n=1}^\infty \dfrac{1}{n}$ 发散, 所以 $\sum\limits_{n=1}^\infty \sin\dfrac{\pi}{n}$ 也发散.
\end{column}
\end{columns}

\end{frame}

%------------------------------------------------

\begin{frame}{比较判别法的极限形式}

$x_n \leqslant Ay_n$ 的一种特殊情况是 $\lim\limits_{n\to\infty} \dfrac{x_n}{y_n}$ 等于某个有限实数, 或是发散到 $+\infty,$ 这也是我们经常能遇到的情况. 比较判别法在这种情况下, 有更精细一些的表达:

\pause

\begin{thm}[比较判别法的极限形式]
设 $\sum\limits_{n=1}^\infty x_n$ 与 $\sum\limits_{n=1}^\infty y_n$ 是两个正项级数, 且
\[\lim\limits_{n\to\infty} \dfrac{x_n}{y_n} = \ell \quad (0 \leqslant \ell \leqslant +\infty)\]
\begin{enumerate}
\item[(1)] 若 $0 \leqslant \ell < +\infty,$ 则当 $\sum\limits_{n=1}^\infty y_n$ 收敛时, $\sum\limits_{n=1}^\infty x_n$ 也收敛;
\item[(2)] 若 $0 < \ell \leqslant +\infty,$ 则当 $\sum\limits_{n=1}^\infty y_n$ 发散时, $\sum\limits_{n=1}^\infty x_n$ 也发散.
\end{enumerate}
所以当 $0 < \ell < +\infty$ 时, $\sum\limits_{n=1}^\infty x_n$ 与 $\sum\limits_{n=1}^\infty y_n$ 有相同的敛散性.
\end{thm}

\end{frame}

%------------------------------------------------

\begin{frame}{比较判别法极限形式的应用}

\begin{eg}
\begin{center}
判断正项级数 $\sum\limits_{n=1}^\infty \left( e^{\frac{1}{n^2}} - \cos\dfrac{\pi}{n} \right)$ 的敛散性.
\end{center}
\end{eg}

\pause

\startproof[bg=yellow, fg=red](解)
对通项有估计
\begin{align*}
e^{\frac{1}{n^2}} - \cos\dfrac{\pi}{n} & = \left[ 1 + \dfrac{1}{n^2} + \mathrm{o}\left(\dfrac{1}{n^2}\right) \right] - \left[ 1 - \dfrac{1}{2} \left(\dfrac{\pi}{n}\right)^2 + \mathrm{o}\left(\dfrac{1}{n^2}\right) \right] \\
& = \left( 1 + \dfrac{\pi^2}{2} \right) \dfrac{1}{n^2} + \mathrm{o}\left(\dfrac{1}{n^2}\right),
\end{align*}
\pause
所以
\[\lim\limits_{n\to\infty} \dfrac{e^{\frac{1}{n^2}} - \cos\dfrac{\pi}{n}}{\dfrac{1}{n^2}} = 1 + \dfrac{\pi^2}{2}.\]
依据比较判别法的极限形式, 由 $\sum\limits_{n=1}^\infty \dfrac{1}{n^2}$ 收敛知 $\sum\limits_{n=1}^\infty \left( e^{\frac{1}{n^2}} - \cos\dfrac{\pi}{n} \right)$ 也收敛.
\qed

\end{frame}

%------------------------------------------------

\section[积分判别法]{正项级数敛散性的积分判别法}

%------------------------------------------------

\begin{frame}{正项级数敛散性的积分判别法}

\only<1>{
无穷级数与我们上一章学习的反常积分有着密切的联系. 我们接下来要介绍一种判别法, 建立起二者敛散性分析的桥梁. 利用它, 我们可以通过反常积分的敛散性判别正项级数的敛散性, 反之亦然.

\vspace{1em}

我们很快会看到, 前面的很多例子中关于正项级数敛散性, 可以很快地通过某些反常积分的敛散性进行判断.
}

\only<2->{
\begin{thm}[积分判别法]
设 $f(x)$ 定义于 $[a, +\infty),$ 且 $f(x) \geqslant 0,$ 进一步假设 $f(x)$ 在任意有限区间 $[a, A]$ 上 Riemann 可积. 设 $\{a_n\}$ 是单调增加且趋于 $+\infty$ 的数列:
\[a = a_1 < a_2 < a_3 < \cdots < a_n < \cdots,\]
令 $\displaystyle u_n = \int_{a_n}^{a_{n+1}} f(x) \mathrm{d} x.$ 那么反常积分 $\displaystyle \int_a^{+\infty} f(x) \mathrm{d} x$ 与正项级数 $\sum\limits_{n=1}^\infty u_n$ 同时收敛或同时发散到 $+\infty,$ 且有
\[\int_a^{+\infty} f(x) \mathrm{d} x = \sum\limits_{n=1}^\infty u_n = \sum\limits_{n=1}^\infty \int_{a_n}^{a_{n+1}} f(x) \mathrm{d} x.\]
\end{thm}
}

\only<3->{
特别地, 当 $f(x)$ 在 $[a, +\infty)$ 上单调减小, 且取 $a_n = n,$ 那么反常积分 $\displaystyle \int_a^{+\infty} f(x) \mathrm{d} x$ 与正项级数 $\sum\limits_{n=[a]+1}^\infty f(n)$ 具有相同的敛散性.
}

\end{frame}

%------------------------------------------------

\begin{frame}{积分判别法的证明}
设正项级数 $\sum\limits_{n=1}^\infty u_n$ 的部分和数列为 $\{S_n\},$ 那么对任意 $A > a,$ 存在 (由 $A$ 唯一确定的) 正整数 $n$ 使得 $a_n \leqslant A < a_{n+1},$ 于是有
\[S_{n-1} \leqslant \int_a^A f(x) \mathrm{d} x \leqslant S_n.\]
\pause
若正项级数 $\sum\limits_{n=1}^\infty u_n$ 收敛, 即数列 $\{S_n\}$ 收敛时, 根据极限的夹逼性, $\displaystyle \int_a^A f(x) \mathrm{d} x$ 当 $A \to +\infty$ 时也收敛, 且收敛到同一极限. \\
若 $\sum\limits_{n=1}^\infty u_n$ 发散到 $+\infty,$ 由不等式 $S_{n-1} \leqslant \int_a^A f(x) \mathrm{d} x$ 可知反常积分 $\displaystyle \int_a^{+\infty} f(x) \mathrm{d} x = \lim\limits_{A\to+\infty} \int_a^{A} f(x) \mathrm{d} x = +\infty.$

\end{frame}

%------------------------------------------------

\begin{frame}{积分判别法的证明}

当 $f(x)$ 在 $[a, +\infty)$ 上单调减小, 且取 $a_n = n$ 时, 在闭区间 $[n, n+1]$ 上有
\[f(n+1) \leqslant f(x) \leqslant f(n),\]
由定积分的保序性, 有不等式
\[f(n+1) = \int_{n}^{n+1} f(n+1) \mathrm{d} x \leqslant \int_{n}^{n+1} f(x) \mathrm{d} x \leqslant \int_{n}^{n+1} f(n) \mathrm{d} x = f(n).\]
\pause
由比较判别法 (使用两次) 知 $\sum\limits_{n=[a]+1}^\infty f(n)$ 与 $\sum\limits_{n=[a]+1}^\infty u_n$ 同时收敛或同时发散, 从而与反常积分 $\displaystyle \int_a^{+\infty} f(x) \mathrm{d} x$ 有相同的敛散性.
\qed

\end{frame}

%------------------------------------------------

\begin{frame}{积分判别法的应用}

我们先把之前用比较判别法做过的例题重新审视一遍.

\pause
\vspace*{0.6em}

\begin{eg}
\begin{center}
判断正项级数 $\sum\limits_{n=1}^\infty \sin\dfrac{\pi}{n}$ 的敛散性.
\end{center}
\end{eg}

\pause
\vspace{0.6em}

\startproof[bg=yellow, fg=red](解)
由于 $f(x) = \sin\dfrac{\pi}{x}$ 在 $[2, +\infty]$ 上非负单减, 所以 $\displaystyle \int_2^{+\infty} f(x) \mathrm{d} x$ 与 $\sum\limits_{n=2}^\infty \sin\dfrac{\pi}{n}$ 有相同的敛散性. 而
\begin{align*}
\int_2^{+\infty} f(x) \mathrm{d} x & = \int_2^{+\infty} \sin\dfrac{\pi}{x}~\mathrm{d} x = \int_{1/2}^{0} \sin\pi t~\mathrm{d}\dfrac{1}{t} = \int_0^{1/2} \dfrac{\sin\pi t}{t^2} \mathrm{d} t \\
& \geqslant \int_0^{1/2} \dfrac{\frac{2}{\pi}t}{t^2} \mathrm{d} t = \dfrac{2}{\pi} {\color{red} \int_0^{1/2} \dfrac{1}{t} \mathrm{d} t} ~~ \text{发散}
\end{align*}
\qed

\end{frame}

%------------------------------------------------

\begin{frame}{积分判别法的应用}

\begin{eg}
\begin{center}
判断正项级数 $\sum\limits_{n=1}^\infty \left( e^{\frac{1}{n^2}} - \cos\dfrac{\pi}{n} \right)$ 的敛散性.
\end{center}
\end{eg}

\pause
\vspace{0.6em}

\startproof[bg=yellow, fg=red](解)
在 $[0, 1]$ 区间上, $g(x) = e^{x^2} - \cos\pi x$ 是非负递增的, 因为很容易算得 $g(0) = 0,$ 且 $g'(x) = 2xe^{x^2} + \pi \sin\pi x \geqslant 0.$ 于是
\[f(x) = g(1/x) = e^{\frac{1}{x^2}} - \cos\dfrac{\pi}{x}\]
为 $[1, +\infty]$ 上单调递减的非负函数. 那么根据积分判别法, 我们约化到了判断反常积分 $\displaystyle \int_1^{+\infty} f(x) \mathrm{d} x$ 的敛散性:
\[
\int_1^{+\infty} f(x) \mathrm{d} x = \int_1^{+\infty} \left( e^{\frac{1}{x^2}} - \cos\dfrac{\pi}{x} \right) \mathrm{d} x = \int_0^1 \dfrac{e^{t^2} - \cos\pi t}{t^2} \mathrm{d} t
\]

\end{frame}

%------------------------------------------------

\begin{frame}{积分判别法的应用}

在可能得奇点 $0$ 附近, 有无穷小量的等价
\begin{align*}
e^{t^2} - \cos\pi t & = (e^{t^2} - 1) + (1 - \cos\pi t) = (e^{t^2} - 1) + 2 \sin^2 \dfrac{\pi t}{2} \\
& \sim t^2 + 2 \cdot \left( \dfrac{\pi t}{2} \right)^2 = \left( 1 + \dfrac{\pi^2}{2} \right) t^2,
\end{align*}
即在 $0$ 附近成立
\[\dfrac{e^{t^2} - \cos\pi t}{t^2} \sim 1 + \dfrac{\pi^2}{2},\]
于是原来的反常积分收敛, 进而知级数 $\sum\limits_{n=1}^\infty \left( e^{\frac{1}{n^2}} - \cos\dfrac{\pi}{n} \right)$ 收敛.
\qed

\end{frame}

%------------------------------------------------

\begin{frame}{积分判别法的应用}

目前我们都是利用积分判别法, 从反常积分的敛散性来判别相应正项级数的敛散性. 事实上, 我们也可以反其道而行之, 由已知敛散性的级数出发, 去判断某些反常积分的敛散性.

\pause
\vspace{0.6em}

\begin{eg}
\begin{center}
证明反常积分 $\displaystyle \int_0^{+\infty} \dfrac{\mathrm{d} x}{1 + x^2\sin^2 x}$ 发散.
\end{center}
\end{eg}

\pause
\vspace{0.6em}

\startproof[bg=yellow, fg=red]
取 $a_n = n\pi, ~ n = 0, 1, 2, \dots,$ 那么
\begin{align*}
u_n & := \int_{n\pi}^{(n+1)\pi} \dfrac{\mathrm{d} x}{1 + x^2\sin^2 x} = \int_{0}^{\pi} \dfrac{\mathrm{d} t}{1 + (n\pi + t)^2\sin^2 t} \\
& > \int_{0}^{{\color{red} \frac{1}{(n+1)\pi}}} \dfrac{\mathrm{d} t}{1 + (n\pi + t)^2\sin^2 t}
\end{align*}
\end{frame}

%------------------------------------------------

\begin{frame}{积分判别法的应用}

当 $0 < t < \dfrac{1}{(n+1)\pi} < 1$ 时, 如下不等式恒成立:
\[(n\pi + t)^2 \sin^2 t < (n + 1)^2 (\pi t)^2 < (n + 1)^2 \pi^2 \dfrac{1}{(n+1)^2\pi^2} = 1.\]
于是
\[u_n > \int_{0}^{\frac{1}{(n+1)\pi}} \dfrac{\mathrm{d} t}{1 + (n\pi + t)^2\sin^2 t} > \dfrac{1}{2\pi} \cdot \dfrac{1}{n+1}.\]
由于 $\sum\limits_{n=1}^\infty \dfrac{1}{n+1}$ 发散, 有比较判别法知 $\sum\limits_{n=1}^\infty u_n$ 发散,又由积分判别法知, 反常积分 $\displaystyle \int_0^{+\infty} \dfrac{\mathrm{d} x}{1 + x^2\sin^2 x}$ 发散.
\qed

\pause

\begin{question}[思考题]
\begin{center}
证明反常积分 $\displaystyle \int_0^{+\infty} \dfrac{\mathrm{d} x}{1 + x^{\color{red}4}\sin^2 x}$ 收敛.
\end{center}
\end{question}

\end{frame}

%------------------------------------------------

\section{总结}

%------------------------------------------------

\begin{frame}{总结}

\begin{itemize}
\item[{\larger[2] \color{red} \ding{43}}] 正项级数的定义: 通项都是非负实数的级数.
\pause
\vspace{0.8em}
\item[{\larger[2] \color{red} \ding{43}}] 正项级数的收敛原理: 正项级数收敛的充分必要条件是它的部分和数列有上界.
\pause
\vspace{0.8em}
\item[{\larger[2] \color{red} \ding{43}}] 不等式形式的比较判别法、极限形式的比较判别法: 通过通项之间的比较, 由已知敛散性或者敛散性容易判断的级数的敛散性, 导出我们感兴趣的级数的敛散性.
\pause
\vspace{0.8em}
\item[{\larger[2] \color{red} \ding{43}}] 积分判别法: 建立了非负函数反常积分敛散性与正项级数敛散性之间的等价关系.
\pause
特别地, 当函数有单调性 (单调递减) 时, 有
\begin{empheq}[box={\fancymath[colback=cyan!30, drop lifted shadow, sharp corners]}]{equation*}
\color{red!80} \int_a^{+\infty} f(x) \mathrm{d} x ~~ \text{收敛} ~~ \Longleftrightarrow ~~ \sum\limits_{n=[a]+1} f(n) ~~ \text{收敛}
\end{empheq}
\end{itemize}

\end{frame}

%------------------------------------------------

%------------------------------------------------
% % closing page
% \begin{frame}
% \HUGE{\centerline{\bfseries {\color{gray} The} {\color{zkblue} End}}}

% \pause

% \vspace{1.2em}

% \Large{\centerline{\bfseries {\color{gray}Thank} \quad {\color{zkblue} You}}}

% \end{frame}

%----------------------------------------------------------------------------------------

\end{document}
